\documentclass[12pt,a4paper]{exam}
\usepackage[T1]{fontenc}
\usepackage{longtable}
\usepackage{graphicx}

\usepackage{hyperref}
%\urlstyle{same}  % don't use monospace font for urls
\setlength{\parindent}{0pt}
\setlength{\parskip}{6pt plus 2pt minus 1pt}

\usepackage{geometry}
\usepackage{multicol}
\usepackage{relsize}
\usepackage{pgf}
\usepackage{minted}

\geometry{
  %includeheadfoot,
  margin=2.54cm
}
\title{Exercises on Design Patterns IV}
\author{Nearly there --- just a few to go\ldots}
\date{2018}

\fvset{tabsize=2}
%\fvset{fontsize=\relsize{-1}}
%\fvset{frame=single}

% https://www.javacodegeeks.com/2015/09/decorator-design-pattern.html

\newcommand{\code}{/Users/keith/Documents/Courses/sdp/repo/worksheets/dp-four/kotlin/src}
%\newcommand{\codesol}{/Users/keith/Courses/sdp/SDP-2017/exercises/week11sol/src/main/scala}

\begin{document}
\maketitle

In these exercises we will be examining the following design patterns:
\begin{itemize}
\item Flyweight, 
\item Memento, and
\item Proxy.
\end{itemize}
The source code snippets listed below are written using the Kotlin language; 
similar versions in Scala can be found on the source 
code repository under the \verb!worksheets/dp-four! folder. 

\begin{questions}

\question
This question concerns the \textsc{Flyweight} design pattern.

Sometimes the objects in an application might have great similarities and be of a similar kind 
(a similar kind here means that most of their properties have similar values and only a few of 
them vary in value). In case they are also heavy objects to create, they should be controlled by 
the application developer. Otherwise, they might consume much of the memory and eventually slow 
down the whole application.

The \textsc{Flyweight Pattern} is designed to control such kind of object creation and 
provides you 
with a basic caching mechanism. It allows you to create one object per type (the type here 
differs by a property of that object), and if you ask for an object with the same property 
(already created), it will return you the same object instead of creating a new one.

The X-programming site allows users to create and execute programs using their favourite
programming language. It provides you with plenty of programming language options. You choose 
one, write a program with it and execute it to see the result.

But now the site has started losing its users, the reason being the slowness of the site. The 
users are not interested in it any more. The site is very popular and sometimes there could be 
more than thousands of programmers using it. Because of that, the site is crawling. But the heavy 
usage is not the real problem behind the slowness of the site. Let us see the core programming of 
the site which allows users to run and execute their program, and the true issue will be revealed 
there.
\VerbatimInput{\code/flyweight/Code.kt}
The above class is used to set the code done by the programmer in order to get it executed. The 
\mintinline{scala}!Code! object is a lightweight simple object having a property code along 
with its setter and getter.
\VerbatimInput{\code/flyweight/Platform.kt}
The \mintinline{scala}!Platform! interface is implemented by the language specific 
platform in order to execute the code. It has one method, executes, which takes 
the \mintinline{scala}!Code! object as its parameter.
\VerbatimInput{\code/flyweight/ScalaPlatform.kt}
The above class implements the \mintinline{scala}!Platform! 
interface and provides an implementation for the \mintinline{scala}!execute! 
method, to execute the code in Scala.

To execute the code, a \mintinline{scala}!Code! object which contains the code and a 
\mintinline{scala}!Platform! object to execute the code get created. The code looks like this:
\begin{Verbatim}
Platform platform = ScalaPlatform()
platform.execute(code)
\end{Verbatim}
The intent of the \textsc{Flyweight Pattern} is to use shared objects to support large numbers 
of fine-grained objects efficiently. A flyweight is a shared object that can be used in 
multiple contexts simultaneously. The flyweight acts as an independent object in each context 
--- it’s indistinguishable from an instance of the object that’s not shared. 
Flyweights cannot make assumptions about the context in which they operate. 

The key concept here is the distinction between \emph{intrinsic} and 
\emph{extrinsic} state. 
\begin{description}
\item[Intrinsic state] is stored in the flyweight; it consists of information that’s 
independent of the flyweight’s context, thereby making it sharable. 
\item[Extrinsic state] depends on and varies with the flyweight’s context and therefore can’t be 
shared. 
\end{description}
Client objects are responsible for passing extrinsic state to the flyweight when it needs it.

To solve the problem, you will need to provide a platform factory class which will control 
the creation of the \mintinline{scala}!Platform! objects.
\VerbatimInput{\code/flyweight/PlatformFactory.kt}
Your code should pass the following (indicative) tests:
\VerbatimInput{\code/flyweight/TestFlyweight.kt}
\begin{solution}
	
\end{solution}

\question
This question concerns the \textsc{Memento} design pattern.

Sometimes it’s necessary to record the internal state of an object. 
This is required when implementing checkpoints and ``undo'' mechanisms that let users back 
out of tentative operations or recover from errors. You must save state information somewhere, 
so that you can restore objects to their previous conditions. But objects normally 
encapsulate some or all of their state, making it inaccessible to other objects and 
impossible to save externally. Exposing this state would violate encapsulation, 
which can compromise the application’s reliability and extensibility.

The \textsc{Memento} pattern can be used to accomplish this without exposing the 
object’s internal structure. The object whose state needs to be captured is referred to 
as the originator.

You will create a class that will contain two double type fields and we will run 
some mathematical operations on it. We will provide users with an \mintinline{scala}!undo!
operation. If the results after some operations are not satisfactory for a user, 
the user can call the \mintinline{scala}!undo! operation which will restore the state of 
the object to the last saved point.
\VerbatimInput{\code/memento/Originator.kt}
The above is the \mintinline{scala}!Originator! class whose object state should be saved 
in a memento. The class contains two double type’s fields \mintinline{scala}!x! and 
\mintinline{scala}!y!, and also takes a reference of a \mintinline{scala}!CareTaker!. 
The \mintinline{scala}!CareTaker! is used to save and retrieve the memento objects 
which represent the state of the \mintinline{scala}!Originator! object.

In the constructor, you should save the initial state of the object using the 
\mintinline{scala}!createSavepoint! method. This method creates a memento object and 
requests the care taker to take care of the object. You should use a 
\mintinline{scala}!lastUndoSavepoint! variable which is used to store the key name of 
last saved memento in order to implement the \mintinline{scala}!undo! operation.

The class provides three types of \mintinline{scala}!undo! operations:
\begin{itemize}
\item The \mintinline{scala}!undo! method without any parameter restores the last saved state.
\item The \mintinline{scala}!undo! with the \mintinline{scala}!savepoint! name as a parameter
restores the state saved with that particular \mintinline{scala}!savepoint! name. 
\item The \mintinline{scala}!undoAll! method asks the care taker to clear all the 
savepoints and set it to the initial state (the state at the time of the creation of the object).
\end{itemize}
\VerbatimInput{\code/memento/Memento.kt}
The \mintinline{scala}!Memento! class is used to store the state of the 
\mintinline{scala}!Originator! and stored by the care taker. 
The class does not have any setter methods, it is only used to get the state of the object.
\VerbatimInput{\code/memento/CareTaker.kt}
The above class is the care taker class used to store and provide the 
requested memento object. The class contains the \mintinline{scala}!saveMemento! 
method is used to save the memento object, the \mintinline{scala}!getMemento! is 
used to return the request memento object and the \mintinline{scala}!clearSavepoints! 
method which is used to clear all the savepoints and it deletes all the saved memento objects.

Your code should pass the following (indicative) tests:
\VerbatimInput{\code/memento/TestMementoPattern.kt}
The above code will result to the following output:
\begin{Verbatim}
Saving state...INITIAL
Default State: X: 5.0, Y: 10.0
State: X: 510.0, Y: 10.0
Saving state...SAVE1
State: X: 510.0, Y: 23.181818181818183
Undo at ...SAVE1
State after undo: X: 510.0, Y: 10.0
Saving state...SAVE2
State: X: 1.32651E8, Y: 10.0
Saving state...SAVE3
State: X: 1.32651E8, Y: 1.3265097E8
Saving state...SAVE4
State: X: 1.32651E8, Y: 6029590.909090909
Undo at ...SAVE2
Retrieving at: X: 1.32651E8, Y: 10.0
Undo at ...INITIAL
Clearing all save points...
State after undo all: X: 5.0, Y: 10.0	
\end{Verbatim}

\begin{solution}
	
\end{solution}

\question
The \textsc{Proxy} design pattern comes up with several different variations. 
Some of the important variations are:
\begin{itemize}
\item Remote Proxy, 
\item Virtual Proxy, and 
\item Protection Proxy. 
\end{itemize}
In this question, we will examine these variations and you will implement each of them in Scala.
\begin{parts}
\part \textsc{Remote Proxy}\\

There is a Pizza Company, which has its outlets at various locations. 
The owner of the company gets a daily report by the staff members of the company from 
various outlets. The current application supported by the Pizza Company is a desktop 
application, not a web application. So, the owner has to ask his employees to generate the 
report and send it to him. But now the owner wants to generate and check the report by his own, 
so that he can generate it whenever he wants without anyone’s help. The owner wants you to 
develop an application for him.

The problem here is that all applications are running at their respective JVMs and the 
\textsc{Report Checker} application (which you will design) should run in the owner’s local 
system. The object required to generate the report does not exist in the owner’s system 
JVM and you cannot directly call on the remote object. 
You should use \textsc{Remote Proxy} to solve this problem.

The \mintinline{scala}!ReportGenerator! interface should looks like this:
\VerbatimInput{\code/proxy/remoteproxy/ReportGenerator.kt}
This is a remote interface which defines the methods that a client can call remotely. 
It’s what the client will use as the class type for your service. 
Both the stub and actual service will implement this. 
The method in the interface returns a \mintinline{scala}!String! object. 
You can return any object from the method; this object is going to be shipped over the wire 
from the server back to the client, so it must be \mintinline{scala}!Serializable!. 
Please note that all the methods in this interface should throw 
\mintinline{scala}!RemoteException!.
\VerbatimInput{\code/proxy/remoteproxy/ReportGeneratorImpl.kt}
The final step is to start the service that is, run your concrete implementation of 
remote class, in this case the class is \mintinline{scala}!ReportGeneratorImpl!.

So far, you have created and run a service. You now need a client to use it. 
The report application for the owner of the pizza company will use this service 
in order to generate and check the report. 
\VerbatimInput{\code/proxy/remoteproxy/ReportGeneratorClient.kt}
The above program should result in the following output:
\begin{Verbatim}
********************Location X Daily Report********************
 Location ID: 012
 Today's Date: Sun Sep 39 00:11:23 GMT 2016
 Total Pizza Sell: 112
 Total Sale: $2534
 Net Profit: $1985
***************************************************************	
\end{Verbatim}

\part \textsc{Virtual Proxy}\\

The \textsc{Virtual Proxy} design pattern is a memory saving technique that recommends 
postponing an object creation until it is needed; it is used when creating an 
object the is expensive in terms of memory usage or processing involved. 

Suppose there is a \mintinline{scala}!Company! object in your application and this 
object contains a list of employees of the company in a 
\mintinline{scala}!ContactList! object. 
There could be thousands of employees in a company. 
Loading the \mintinline{scala}!Company! object from the database along with the list of all 
its employees in 
the \mintinline{scala}!ContactList! object could be very time consuming. 
In some cases you don’t even require the list of the employees, but you are forced to wait 
until the company and its list of employees loaded into the memory.

One way to save time and memory is to avoid loading of the employee objects until required, 
and this is done using the \textsc{Virtual Proxy}. 
This technique is also known as \textsc{Lazy Loading} where you are fetching the 
data only when it is required.
\VerbatimInput{\code/proxy/virtualproxy/Company.kt}
The above \mintinline{scala}!Company! class has a reference to 
\mintinline{scala}!ContactList! interface whose real object will be load only when requested 
to call of \mintinline{scala}!contactList()! method.
\VerbatimInput{\code/proxy/virtualproxy/ContactList.kt}
The \mintinline{scala}!ContactList! interface only contains one method 
\mintinline{scala}!getEmployeeList()! which is used to get the employee list of 
the company.
\VerbatimInput{\code/proxy/virtualproxy/ContactListImpl.kt}
The above class should create a real \mintinline{scala}!ContactList! object which will 
return the list of employees of the company. 
The object will be loaded on demand, only when required.
\VerbatimInput{\code/proxy/virtualproxy/ContactListProxyImpl.kt}
The \mintinline{scala}!ContactListProxyImpl! is the proxy class which also implements 
\mintinline{scala}!ContactList! and holds a reference to the real 
\mintinline{scala}!ContactList! object. On the implementation of the method 
\mintinline{scala}!getEmployeeList()! it will check if the 
\mintinline{scala}!contactList! reference is \mintinline{scala}!null!, 
then it will create a real \mintinline{scala}!ContactList! object and then will 
invoke the \mintinline{scala}!getEmployeeList()! method on it to get the list of the employees.

The \mintinline{scala}!Employee! class should look like this, although you will need to 
complete the method.
\VerbatimInput{\code/proxy/virtualproxy/Employee.kt}
You should then test your code with:
\VerbatimInput{\code/proxy/virtualproxy/TestVirtualProxy.kt}
Producing the following output:
{\small\begin{Verbatim}
Company object created...
Company Name: ABC Company
Company Address: Alabama
Company Contact No.: 011-2845-8965
Requesting for contact list
Creating contact list and fetching list of employees...
Employee Name: Employee A, EmployeeDesignation: SE, Employee Salary: 2565.55
Employee Name: Employee B, EmployeeDesignation: Manager, Employee Salary: 22574.0
Employee Name: Employee C, EmployeeDesignation: SSE, Employee Salary: 3256.77
Employee Name: Employee D, EmployeeDesignation: SSE, Employee Salary: 4875.54
Employee Name: Employee E, EmployeeDesignation: SE, Employee Salary: 2847.01	
\end{Verbatim}
}

\part \textsc{Protection Proxy}\\

In general, objects in an application interact with each other to implement the 
overall application functionality. Most application object are generally accessible to 
all other objects in the application. At times, it may be necessary to restrict 
the accessibility of an object only to a limited set of client objects based on their 
access rights. When a client object tries to access such an object, 
the client is given access to the services provided by the object only if the client can 
furnish proper authentication credentials. 
In such cases, a separate object can be designated with the responsibility of verifying 
the access privileges of different client objects when they access the actual object. 
In other words, every client must successfully authenticate with this designated object to 
get access to the actual object functionality. Such an object with which a client needs 
to authenticate to get access to the actual object can be referred as an object 
authenticator which is implemented using the \textsc{Protection Proxy} design pattern.

We will return to the \mintinline{scala}!ReportGenerator! application that you worked on
earlier for the pizza company. The owner now requires that only he can generate the daily report 
and that no other employee should be able to do so.

To implement this security feature, you are going to use the \textsc{Protection Proxy} 
design pattern which will check if the object which is trying to generate the report is 
the owner; in this case, the report gets generated, otherwise it will not.
\VerbatimInput{\code/proxy/protectedproxy/Staff.kt}
The \mintinline{scala}!Staff! interface is used by the 
\mintinline{scala}!Owner! and the \mintinline{scala}!Employee! 
classes and the interface has two methods: \mintinline{scala}!isOwner()! returns a 
boolean to check whether the calling object is the owner or not. 

The other method is used to set the \mintinline{scala}!ReportGeneratorProxy! 
which is a protection proxy used to generate the report is \mintinline{scala}!isOwner()! 
method return true.
\VerbatimInput{\code/proxy/protectedproxy/Employee.kt}
The \mintinline{scala}!Employee! class implements the \mintinline{scala}!Staff! interface, 
since it's an employee \mintinline{scala}!isOwner()! method return \mintinline{scala}!false!. 
The \mintinline{scala}!generateDailyReport()! method asks 
\mintinline{scala}!ReportGenerator! to generate the daily report.
\VerbatimInput{\code/proxy/protectedproxy/Owner.kt}
The \mintinline{scala}!Owner! class looks almost same as the 
\mintinline{scala}!Employee! class, the only difference is that the 
\mintinline{scala}!isOwner()! method returns \mintinline{scala}!true!.
\VerbatimInput{\code/proxy/protectedproxy/ReportGeneratorProxy.kt}
The \mintinline{scala}!ReportGeneratorProxy! acts as a 
\textsc{Protection Proxy} which has only one method \mintinline{scala}!generateDailyReport()!
that is used to generate the report.
\VerbatimInput{\code/proxy/protectedproxy/ReportGeneratorProtectionProxy.kt}
The above class is the concrete implementation of the 
\mintinline{scala}!ReportGeneratorProxy! which holds a reference to the 
\mintinline{scala}!ReportGenerator! interface which is the remote proxy. 
In the \mintinline{scala}!generateDailyReport()! method, it checks if the staff is 
referring to the owner, then asks the remote proxy to generate the report, 
otherwise it returns a string with \verb!“Not Authorised to view report”! as a message.
\VerbatimInput{\code/proxy/protectedproxy/TestProtectionProxy.kt}
The above program should result in the following output:
\begin{Verbatim}
For owner:
********************Location X Daily Report********************
 Location ID: 012
 Today’s Date: Sun Sep 14 13:28:12 IST 2014
 Total Pizza Sell: 112
 Total Sale: $2534
 Net Profit: $1985
***************************************************************
For employee:
Not Authorized to view report.	
\end{Verbatim}
\end{parts}

\begin{solution}
	
\end{solution}

\end{questions}

\end{document}